\chapter{Conclusioni}
\label{cap:conclusioni}

\section{Che cosa è stato stabilito}

Il lavoro non produce un miglioramento delle metriche. Produce cinque
affermazioni sostenute da misure, che vale la pena enunciare nella loro forma
più compatta.

\subsection{Le metriche su ASLG-PC12 sono sature}

Non nel senso vago di ``i punteggi sono alti'', ma nel senso preciso che una
baseline a regole di poche righe, che non apprende nulla, ottiene ROUGE-L
$0{,}9697$ e supera ogni risultato ottenuto con reinforcement learning
(Capitolo~\ref{cap:dataset}). Persino la trasformazione
$\mathrm{uppercase}(\text{text})$ ottiene $0{,}6454$, cioè batte il miglior
risultato con RL ($0{,}6076$).

La causa è documentata: il corpus è sintetico
(sezione~\ref{sec:sintetico}), la sovrapposizione fra token del gloss e del
testo è $98{,}23\%$ contro il $15{,}59\%$ di un corpus annotato da umani
\cite{peng2023monolingual}, e conoscere il token sorgente rimuove circa il $90\%$
dell'incertezza sul gloss ($0{,}856$ bit residui su $8{,}578$).

La metrica stessa contribuisce: l'implementazione di ROUGE-L impiegata è
insensibile al maiuscolo e spezza i token sui caratteri non alfanumerici, quindi
assegna $1{,}00$ a un'uscita in minuscolo identica al testo di partenza
(sezione~\ref{sec:difetto-rouge}).

Il segno più chiaro della saturazione è che un modello da $0{,}5$ miliardi di
parametri, con adattatori LoRA, addestrato per tre epoche su una singola GPU,
supera il miglior valore pubblicato su questo corpus.

\subsection{La domanda di ricerca iniziale era mal posta}

La domanda ``il GRPO migliora la generazione di gloss?'' non è rispondibile su
ASLG-PC12, perché presuppone che esista un margine di miglioramento su cui il
reinforcement learning possa agire. Con l'SFT a exact match $0{,}8130$ e validity
$0{,}9998$, e una baseline a regole a ROUGE-L $0{,}9685$, quel margine non
esiste.

La domanda si riduce, su questo corpus, a ``il GRPO memorizza una mappa
quasi-identità meglio dell'SFT?'', che è una questione di ottimizzazione e non
riguarda la lingua dei segni.

\subsection{Il risultato negativo ha un meccanismo in tre parti}

Questo è il contributo sostanziale. Non ``il GRPO non ha funzionato'', ma tre
misure che spiegano \emph{perché} non poteva funzionare.

\begin{description}
  \item[Parte 1: il collasso del supporto.] Il gradiente del GRPO è
        proporzionale al vantaggio~\eqref{eq:grpo-adv}, che si annulla quando
        tutti i candidati di un gruppo ricevono la stessa reward. Sul modello
        base senza vincolo di decodifica il $73{,}9\%$ dei gruppi è morto; sulla
        policy SFT l'$84{,}8\%$ dei gruppi ha varianza nulla per la ragione
        opposta, cioè perché tutti i candidati sono corretti
        (sezione~\ref{sec:supporto}). Il GRPO è schiacciato fra due fallimenti
        simmetrici, e la finestra utile su un compito quasi deterministico è
        stretta.
  \item[Parte 2: la degradazione della pipeline.] Applicare il GRPO alla policy
        SFT porta ROUGE-L da $0{,}9752$ a $0{,}5114$ ed exact match da
        $0{,}8130$ a $0{,}0203$. Il meccanismo è che gli aggiornamenti avvengono
        soltanto sul $15{,}2\%$ dei gruppi con varianza, che è per costruzione il
        sottoinsieme in cui la policy era incerta; ottimizzare la reward su un
        campione così selezionato, con la sola penalità KL come freno, sposta la
        policy fuori dalla regione corretta. È sovraottimizzazione del
        surrogato~\cite{gao2022overoptimization}, prevedibile con una singola
        sonda sui rollout prima di addestrare.
  \item[Parte 3: la confusione dei fattori.] Dei circa $0{,}47$ di ROUGE-L che
        separano il punto di partenza peggiore dal miglior risultato con RL,
        $0{,}33$ vengono dal \emph{prompting} (few-shot con retrieval, nessun
        parametro aggiornato) e circa $0{,}14$ sono attribuiti al RL in modo non
        isolato (sezione~\ref{sec:decomposizione}). Il fattore sperimentale più
        efficace misurato in tutto il progetto non è il metodo di addestramento.
\end{description}

Va aggiunto un chiarimento sui limiti di questa affermazione. TRL~$0{,}24{,}0$
non implementa il \emph{dynamic sampling} di DAPO~\cite{yu2025dapo}, che scarta i
gruppi a varianza nulla e ricampiona, né il \emph{clip-higher}; e la variante di
denominatore effettivamente ottimizzata (\texttt{dapo}) è stata selezionata per
omissione e non per scelta (sezione~\ref{sec:denominatori}). Quindi nessuna delle
mitigazioni note contro il collasso del supporto era attiva. Il risultato
negativo riguarda questa configurazione, non il GRPO in generale.

\subsection{Il vincolo simbolico abilita l'apprendimento senza migliorare le metriche}

Il risultato neuro-simbolico, e il più controintuitivo. Il Trie
(Capitolo~\ref{cap:decoding}) fa \emph{crollare} ROUGE-L da $0{,}3648$ a
$0{,}1373$ e contemporaneamente:

\begin{itemize}
  \item alza la non-copy-token accuracy da $0{,}0006$ a $0{,}0432$, un fattore
        circa $70$;
  \item alza la validity da $0{,}1043$ a $0{,}9006$;
  \item porta i gruppi morti dal $73{,}9\%$ al $19{,}9\%$, cioè \emph{crea} il
        supporto su cui il reinforcement learning può operare.
\end{itemize}

L'interpretazione: il vincolo restringe la distribuzione generativa alla regione
in cui la reward è discriminante. Fuori da quella regione i candidati sono tutti
ugualmente cattivi e il vantaggio li annulla tutti.

La conseguenza pratica ha valore generale. In un esperimento di reinforcement
learning su uscite strutturate, il vincolo simbolico non è un accessorio di
qualità: è una precondizione di fattibilità, e va valutato misurando il supporto
dei rollout invece delle metriche finali. Il costo che il vincolo impone sulla
qualità del contenuto è reale e documentato in
letteratura~\cite{reddy2603,zhou2604,ye2504}; la conclusione di questo lavoro è
che su un compito con supporto scarso quel costo va pagato, perché l'alternativa
è un gradiente nullo.

\subsection{L'onestà metodologica come parte del risultato}

Tre ritrattazioni misurate sono documentate in questa relazione, e sono parte del
contributo.

\begin{enumerate}
  \item L'ipotesi ``la reward premia la verbosità'', ragionevole a priori e
        documentata in letteratura~\cite{singhal2023length}, è stata
        \textbf{refutata} due volte: analiticamente
        (l'equazione~\eqref{eq:append-k} è decrescente, con perdita misurata di
        circa $-0{,}13$ per token appeso) ed empiricamente (i rapporti di
        lunghezza delle celle con RL sono $1{,}011$ e $0{,}886$, cioè non c'è
        deriva). Il problema delle celle con RL è la varianza di lunghezza, non
        il bias.
  \item La stima del prior strutturale è stata \textbf{corretta due volte}
        prima di arrivare a $+0{,}0444$ bit, per contaminazione fra insiemi e per
        un backoff mal specificato (sezione~\ref{sec:qualificazione-strutturale}).
        Un guadagno così piccolo può essere prodotto interamente da un errore di
        metodo, e chi legge deve sapere che l'errore è stato cercato.
  \item Le tre componenti di reward basate su Viterbi sono state
        \textbf{rimosse} dopo aver dimostrato che a lunghezza uguale collassano
        algebricamente sullo stesso segnale, e che empiricamente restano entro
        $\pm 0{,}006$ dal valore di controllo
        (sezione~\ref{sec:reward-rimosse}).
\end{enumerate}

Va dichiarato anche un test rosso: \texttt{test\_ssh\_alias\_no\_user\_no\_port}
fallisce, per una dipendenza dalla rete nel test stesso, e il fallimento è stato
verificato presente anche sul codice non modificato
(Capitolo~\ref{cap:codebase}).

\section{Lavoro futuro}

Le direzioni seguenti derivano direttamente dai risultati, in ordine di
priorità.

\subsection{Cambiare corpus}

La prima e più importante. Il compito T2G va valutato su PHOENIX-2014T, dove la
sovrapposizione lessicale fra gloss e testo è $15{,}59\%$ invece di $98{,}23\%$
e i migliori risultati riportati sono intorno a ROUGE-L
$55{,}18$~\cite{ouargani2023t2g}. Lì esiste un margine reale, e le domande sul
reinforcement learning diventano rispondibili.

Il \emph{caveat} va dichiarato in anticipo: PHOENIX-2014T è tedesco e riguarda la
lingua dei segni tedesca. Il cambio di corpus comporta un cambio di lingua
sorgente, quindi la scelta del modello base e l'analisi degli artefatti vanno
rifatte da zero. Non è un trasferimento gratuito.

Una direzione più radicale, coerente con la critica al gloss come
rappresentazione~\cite{yinread2020stmc,yin2021including}, è di abbandonare
l'intermedio testuale. Esistono lavori che applicano il rinforzo alla traduzione
della lingua dei segni senza passare per il gloss~\cite{rao2025rvlf}, usando un
segnale di qualità definito su rappresentazioni visive. In quel contesto il
problema del supporto dei rollout andrebbe riesaminato da zero, perché lo spazio
delle uscite e la reward sono di natura completamente diversa; ma la domanda
sarebbe posta sul compito reale invece che su una sua trascrizione impoverita.

\subsection{Ridisegnare l'esperimento primario}

L'esperimento primario dovrebbe essere una griglia
\begin{equation}
  \text{prompting} \times \text{decodifica}
\end{equation}
con il metodo di addestramento come fattore \emph{secondario}. La ragione è
diretta: sono i due fattori con l'effetto misurato più grande
($+0{,}33$ il prompting, e il vincolo che porta i gruppi morti dal $74\%$ al
$20\%$), e sono gli unici due che il disegno attuale confonde fra loro.

Il metodo di addestramento va applicato \emph{dopo} aver identificato il
substrato con il supporto migliore, che secondo la
sezione~\ref{sec:supporto} è few-shot con Trie (gruppi morti allo $0{,}85\%$,
reward media $-0{,}3233$).

\subsection{Eseguire i gate degli obiettivi ausiliari}

I due obiettivi del Capitolo~\ref{cap:ausiliari} sono implementati, verificati e
non addestrati. I criteri di promozione sono definiti:

\begin{enumerate}
  \item per la massa ammessa, una sonda in inferenza che mostri una riduzione
        misurabile della massa rimossa dal vincolo di decodifica;
  \item per la loss strutturata, un confronto controllato fra testa addestrata
        con transizioni vere e con transizioni mescolate casualmente.
\end{enumerate}

Entrambi i gate costano poco rispetto a un addestramento completo, ed entrambi
possono produrre un risultato negativo che risparmia l'esperimento. Il valore
$+0{,}0444$ bit della qualificazione strutturale suggerisce che il margine
disponibile sia modesto.

\subsection{Il reinforcement learning, se e quando}

L'ordine corretto delle operazioni, ricavato per via negativa da questo lavoro,
è: prima si stabilisce un substrato con supporto misurato, poi si progetta la
reward con test avversari, poi si addestra. Non prima.

La verifica di fattibilità è a costo quasi nullo: si generano $G$ candidati su un
campione di prompt, si calcolano le reward, e si contano i gruppi a varianza
nulla. Se la frazione è alta, il reinforcement learning non apprenderà,
indipendentemente dagli iperparametri. È una misura che richiede minuti e che
avrebbe risparmiato, in questo progetto, una quantità considerevole di tempo di
GPU.

\section{La lezione trasferibile}

Se il lavoro ha un unico insegnamento riutilizzabile, è il seguente.

\begin{quote}
Prima di ottimizzare una metrica, misurare quanto la si può ottenere senza
apprendimento; e prima di addestrare con reinforcement learning, misurare se il
segnale esiste.
\end{quote}

Le due misure hanno costo trascurabile. La prima è una tabella di
corrispondenze contata sul training set, che qui ha rivelato ROUGE-L $0{,}9697$ e
ha reso interpretabile ogni numero successivo. La seconda è una sonda sui
rollout, che qui ha rivelato il $73{,}9\%$ di gruppi morti prima
dell'addestramento e l'$84{,}8\%$ di varianza nulla dopo l'SFT, e che avrebbe
previsto entrambi i risultati negativi.

Nessuna delle due era presente nel disegno originale. Entrambe, in retrospettiva,
erano le prime cose da fare.
